

\documentclass[gmd, manuscript]{copernicus}

\begin{document}


\title{Response letter}

% \Author[affil]{given_name}{surname}

\Author[1,2]{Xiaoxu}{Shi}
\Author[3]{Alexandre}{Cauquoin}
\Author[1,4]{Gerrit}{Lohmann}
\Author[4]{Lukas}{Jonkers}
\Author[1]{Qiang}{Wang}
\Author[1,4]{Hu}{Yang}
\Author[1,5]{Yuchen}{Sun}
\Author[1]{Martin}{Werner}

\affil[1]{Alfred Wegener Institute, Helmholtz Center for Polar and Marine Research, Bremerhaven, Germany}
\affil[2]{Southern Marine Science and Engineering Guangdong Laboratory (Zhuhai), Zhuhai, China}
\affil[3]{Institute of Industrial Science, The University of Tokyo, Kashiwa, Japan}
\affil[4]{MARUM Center for Marine Environmental Sciences, University of Bremen, Bremen, Germany}
\affil[5]{Institute of Tibetan Plateau Research, Chinese Academy of Sciences, Beijing, China}

\runningtitle{MH simulation using AWI-ESM-wiso}

\firstpage{1}

\maketitle


\textbf{Dear Reviewers,}

\textbf{Thank you very much for your positive and constructive comments. In the following, we present our point-to-point responses. Our answers to your comments are written in blue. }

\textbf{Thanks again for your time and efforts.}

 \textbf{Best, }
 
 \textbf{Xiaoxu}

 \medskip

 \section{Comments from Reviewer 1}

Summary of Shi et al.

 \medskip
 
Shi et al. has performed a mid-Holocene (MH) time slice experiment with the new isotope enabled Earth System Model AWI-ESM-wiso. The authors validate the control run and go on to compare the model with isotopic archives. Although the model performs well for the control run and the spatial pattern of the MH anomalies are correct the overall amplitude of the isotope anomalies appear underestimated, which could mean that the MH climate anomalies are underestimated. The authors investigate the spatial and temporal relations of isotopes and climate. One conclusion is that the modern spatial slope can be used as surrogate for the temporal slope in Greenland and Antarctica. Finally, the authors offer an alternative method employing a combination of isotope anomalies and precipitation amount to determine the onset of the West African summer monsoon.

 \medskip
 


 \medskip

General comments.

 \medskip
 
I find the manuscript overall well-written, well-structured and with appropriate figures to illustrate the results. As the detailed comments below show I have a number of suggestions and concerns that I think should be taken into account before considering publication. One major concern is the somewhat superficial treatment of the debate on temporal versus spatial slope, which is a topic that has been extensively researched, and I find some key references missing in this context. Temporal and spatial changes in climate are different in fundamental ways, and this is reflected in the difference in the temporal and spatial slope. I think we are past the point where we are asking if the temporal and spatial slopes are interchangeable, which is summarized by the publications listed below (see comment to L304-307 and L319-324).

In summary I find that minor revisions should be made to accommodate the major and minor comments in this review.

 \textcolor{blue}{\textbf{We really appreciate your positive comments and suggestions, which have significantly improved our manuscript. We have modified our paper accordingly. Details are as follows: }}
 

Detailed comments.

 \medskip

Abstract

 \medskip

1. L2 “Straightforward” do you mean “direct”?

\medskip


\textcolor{blue}{\textbf{Thanks for the comment, we now changed “straightforward” to “direct”}}

%\textcolor{blue}{\textbf{\textit{"However, the simulated responses of the Atlantic Meridional Overturning Circulation (AMOC) in the two models yield different results for both the LIG and the MH."}}}

\medskip

2. L6 Move “well” to right before the comma.
 \medskip
 
\textcolor{blue}{\textbf{We have modified the texts according to the comment.}}

 \medskip

3. L9-10 “The ratio of the MH-PI difference … is reasonable …” clarify if you are comparing to measured d18O.

 \medskip
 
\textcolor{blue}{\textbf{Thanks for the comment, we now changed the text into (see also L10 in the difference-tracked version):}}

\textcolor{blue}{\textbf{\textit{"The ratio of the MH-PI difference in $\delta^{18}$O$_p$ to the MH-PI difference in surface air temperature is comparable to proxy records over Greenland and Antarctica only when summertime air temperature is considered."}}}

 \medskip


4. L39-43 The isotope-temperature relationship can also be affected by the proportion of continental sources and recycling of vapor (Werner et al. (2001), Sjolte et al., (2014)).

 \medskip
 
\textcolor{blue}{\textbf{Thanks for the comment, we now changed the text into (L48-52 in the difference-tracked version):}}

\textcolor{blue}{\textbf{\textit{"Changes in the primary moisture source areas or air mass transport trajectories during the LGM can have a significant impact on the isotope-temperature relationship, which has a potential to invalidate the isotopic paleothermometer approach based on the use of modern observations  \citep{delaygue2000origin,werner2001isotopic}.    Moreover, the isotope-temperature relationship can also be influenced by the proportion of continental sources and recycling of vapor under interglacial boundary conditions \citep{sjolte2014modelling}."}}}

 \medskip


5. L59 “Straightforward” do you mean “direct”?

\medskip

\textcolor{blue}{\textbf{Thanks for the comment, we now changed “straightforward” to “direct”}}

\medskip

6. L139 Is spring equinox fixed to March 21? This information can be added to Table 1.

\medskip

\textcolor{blue}{\textbf{We now added this information into Table 1 (Page 6 in the difference-tracked version).}}

\medskip


7. L172 “Straightforward” do you mean “direct”?

\medskip

\textcolor{blue}{\textbf{In the revised version, we changed “straightforward” to “direct”}}

\medskip

8. Section 2.3.5 “Assimilation product”. I suggest a more informative title such as “Marine d18O reanalysis data”.

\medskip

\textcolor{blue}{\textbf{We have modified the texts according to the comment (L210 in the difference-tracked version). Moreover, the same change was also made for several other parts of the manuscript where “assimilation product” was used. }}

\medskip

9. L214-215 High bias in d18O over Antarctica has also been linked to a low bias in water vapor for cold regions leading to a precipitation weighting towards higher d18O (Masson-Delmotte et al., 2008).
\medskip

\textcolor{blue}{\textbf{Thanks for the comment, we now changed the text into (L234-235 in the difference-tracked version)}}

\textcolor{blue}{\textbf{\textit{"According to \citet{masson2008review}, this  underestimation can also be related to a low bias in water vapor for cold regions that favors higher  $\delta^{18}$O$_p$ values."}}}

\medskip


10. Figure 3/4/5/6. Are the simulated MH-PI  anomalies significant?

\medskip

\textcolor{blue}{\textbf{In the updated manuscript, we performed Student's t-est to test the significance level of the anomalies shown in Fig. 3/4/5/6, and we  marked the area with significance level of greater than 95\%  by black dots. We refer to   Fig. 3/4/5/6 in our revised manuscript and the difference-tracked version.}}

\medskip

11. Figure 6b Why not plot the RMSE between the different species which you write in the text? I think the reader expects this since it is included in the other plots. It doesn’t take much to explain the details in the figure caption.

\medskip

\textcolor{blue}{\textbf{Thanks, we now added the RMSE across all species in Fig. 6b, and we also display the RMSE value for each individual foraminifera type in the corresponding legend. We refer to  Fig. 6b in our revised manuscript  and the difference-tracked version.}}

\medskip


12. L304-307 I think it would be good to move this to the introduction and include some relevant paper discussing this (Cuffey et al., 1992; Sime et al., 2009; Sjolte et al., 2011; Kindler et al., 2014; Sjolte et al., 2014; Guan et al., 2016).

\medskip

\textcolor{blue}{\textbf{Yes, we now added in the introduction (L38-46 in the difference-tracked version):}}

\textcolor{blue}{\textbf{\textit{"The observed present-day spatial slope between isotope and temperature is widely used as a surrogate  for the temporal gradient at reconstruction sites \citep{cuffey1992toward,sime2009evidence,kindler2014temperature}. A model simulation using isoCAM3 suggests that both the temporal slope and spatial slope remain largely stable throughout the last deglaciation  \citep{guan2016understanding}. However, a significant region-dependency is found for both the temporal slope \citep{guan2016understanding} and the spatial slope \citep{sjolte2011modeling}. \citet{guan2016understanding} also point out that the temporal slope is usually smaller than the spatial slope in the extratropics. Besides, some studies indicate that the temporal relationship between isotopes and temperature may vary over time and the temporal isotope-temperature slope may differ from the observed spatial slope at present-day \citep[e.g.,][]{werner2000borehole,sjolte2014modelling}. Therefore, using isotope data from different archives to draw quantitative inferences about past climate variability remains challenging."}}}

\medskip
 

13. Figure 8: It is risky to calculate a slope between only two points, but one could use the STD of the interannual variability to estimate the uncertainty or use the annual/seasonal data to make bootstrap estimates of the uncertainty in slope.

\medskip

\textcolor{blue}{\textbf{In the revised manuscript, we have more analysis on the temporal slopes, we  calculated the slopes not only for annual mean temperature, but also for JJAS and DJFM seasons. For more details please refer to our responses to your comments number 16 and 18.}}

\medskip

14. L323 The estimate is affected by the selection of the data. For example, if coastal data is included or only higher altitude sites on the ice sheet (Sjolte et al., 2011).

\medskip

\textcolor{blue}{\textbf{Thanks for the comment, we now changed the text into (L359-360 in the difference-tracked version):}}

\textcolor{blue}{\textbf{\textit{"In addition,  the spatial $\delta^{18}$O$_p$-temperature relationship  can be affected by the criteria of data selection, e.g., if coastal areas or only higher altitude sites on the ice sheet are included \citep{sjolte2011modeling}."}}}

\medskip


15. L319-324 Vinther et al. (2009) estimated a slope of 0.5 permil/degree C using borehole temperature and ice core data. I think it is well documented from the papers I mention above that the spatial slope and temporal slope are not interchangeable and that there are large regional differences in both. To quote Guan et al. (2016): “Finally, the relation between temporal and spatial slopes is understood using a semiempirical equation that is derived based on both the Rayleigh distillation and a fixed spatial slope. The slope equation quantifies the Boyle's mechanism and suggests that the temporal slope is usually smaller than the spatial slope in the extratropics mainly because of the polar amplification feature in global climate change, such that the response in local temperature at middle and high latitudes is usually greater than that in the total equivalent source temperature.”

\medskip

\textcolor{blue}{\textbf{Thanks, we added this estimation of greenland spatial slope in our text (see also L345-351 in the difference-tracked version):}}

\textcolor{blue}{\textbf{\textit{"Over Greenland, our modeled  $\delta^{18}$O$_p$-temperature gradients under present-day and MH conditions are 0.76$\pm$0.042 and 0.74$\pm$0.045 ‰/$^{\circ}C$, respectively (Fig.7b,g), higher than the  value obtained from modern observations  (0.67 ‰/$^{\circ}C$) \citep{johnsen1989origin} and previous model studies using MPI-ESM-wiso (0.71 ‰/$^{\circ}C$) \citep{cauquoin2019water} and ECHAM4  (0.58 ‰/$^{\circ}C$)  \citep{werner2000borehole}.  However, a much smaller Holocene spatial isotope-temperature slope, ranging from 0.43 ‰/$^{\circ}C$ to 0.53 ‰/$^{\circ}C$, was estimated based on ice core records and borehole temperatures for the Greenland ice sheet  \citep{vinther2009holocene}. "}}}

\medskip

\textcolor{blue}{\textbf{Besides, in the end of section 5.1, as also added (see L361-364 in the difference-tracked version):}}

\textcolor{blue}{\textbf{\textit{"It is important to note that there is a significant region-dependency found for the spatial isotope-temperature slope \citep{sjolte2011modeling}, and that the spatial  slope might differ from the temporal slope due to changes in the seasonality of precipitation   \citep[e.g.,][]{werner2000borehole},  moisture source regions  \citep[e.g.,][]{delaygue2000origin} or  polar amplification feature \citep{guan2016understanding}. In the following section, we aim to analyze the temporal relationship between our simulated $\delta^{18}$O$_p$ and climate variables."}}}

\medskip





16. L341 What is the choice of 0.5 degree C anomalies based on? Can’t you base it on whether the anomalies are significant?

\medskip

\textcolor{blue}{\textbf{Regarding the first question, for numerical reasons, a threshold of 0.5 C (or 1 C) is chosen here to avoid the situation that a small temperature anomaly causes an enormous unrealistic isotope-temperature gradient. On the other hand, the threshold can guarantee that the isotope anomaly is mostly affected by temperature changes rather than other elements such as precipitation. We now clarified this point in the revised paper (see L371-374 in the difference-tracked version): }}

\textcolor{blue}{\textbf{\textit{"To avoid numerical errors in calculated temporal relationships caused by very small MH-PI temperature changes, another criterion is adopted that the absolute change of T between the two time intervals  (i.e., $T_{MH}-T_{PI}$) must be non-negligible, namely not less than 0.5 $^{\circ}C$. Applying such a threshold can also ensure a temperature-dependency of the isotope changes"}}}

\medskip



\renewcommand{\thefigure}{R\arabic{figure}}



    \begin{figure}
      \centering\includegraphics[width=0.75\textwidth,angle=90]{temporal_atm_sig.pdf}
        \caption{(a) Simulated temporal MH-to-PI $\delta^{18}$O$_p$-temperature gradient for all grid boxes with mean annual temperature (MAT) for both PI and MH being lower than 20 $^{\circ}C$ and the change in temperature between MH and PI is significant based on Student's t-test.   (b) Same as in (a), but using the temperature of the warmest month (MTWA). (c) Same as in (a) but using the temperature of the monsoon months (JJAS for the Northern Hemisphere, DJFM for the Southern Hemisphere). (d) Same as in (c) but only for the areas with a temperature anomaly being larger than 0.5 °C.  (e) Same as in (c) but only for the areas with a temperature anomaly being larger than 1 °C Units: ‰/$^{\circ}C$. \label{temporal}}
    \end{figure}

    \textcolor{blue}{\textbf{Regarding the second question, it is a good idea to calculate the isotope-temperature gradient based on the significance of the temperature anomalies. The results are shown here in Fig.~\ref{temporal} in this response letter. From Fig.~\ref{temporal}a we see that the data on Greenland is missing, this is because of the minor change in annual mean temperature between MH and PI, which is statistically not significant for those regions (This can also be seen from the temperature anomaly plot, i.e., Fig. 3c of the paper). But the value for the West Antarctica (0.61 ‰/$^{\circ}C$) are more close to observation  compared to other T definitions. If we only consider the warmest month, as seen in  Fig.~\ref{temporal}b, the pattern is almost the same as in Fig. 8c of the paper, and the detailed numbers for the isotope-temperature gradients are also very similar (Greenland: 0.64 ‰/$^{\circ}C$; Antarctic: 0.48 ‰/$^{\circ}C$; East Antarctica: 0.55 ‰/$^{\circ}C$; West Antarctica: 0.39 ‰/$^{\circ}C$).}}
    
\medskip

17. L363 Judging by Figure 8d there is still quite some regional differences in slope. What is the slope of key ice core locations in Greenland and Antarctica, and what is the reconstructed temperature anomalies from ice core d18O based in this slope? This can then be compared to the simulated temperature anomaly.

\medskip

\textcolor{blue}{\textbf{Yes there are regional difference and it is a good idea to compared modeled temperature anomalies with ice core records. In addition, the slope of key ice core locations in Greenland and Antarctica has been illustrated in Table 2 of the manuscript, and we can use these slopes, along with isotope anomalies documented by ice cores, to compute the temperature anomalies for each ice core location. Please see Fig.~\ref{tp} for detailed information, this figure is also provided in the revised paper (see Fig. 9 of the paper). We also add a table in the paper to illustrate observed temperature anomalies and modeled T anomalies for each region and for each T definition (please refer to Table 3 in the updated paper). In addition, we added the following texts (see also 420-428 in the difference-tracked version):}}


    \begin{figure}
      \centering\includegraphics[width=0.75\textwidth,angle=90]{temp2_anm_proxy.pdf}
        \caption{Simulated  MH-PI T anomalies for (a-c) Greenland and (d-f) Antarctica, with T being (a,d) MAT, (b,e) MTWA and (c.f) T-monsoon.  The marked area has a significance level of greater than 95\% based on Student's t-test. Circles represent reconstructed temperature anomalies based on ice core records. Units: $^{\circ}C$. \label{tp}}
    \end{figure}

\textcolor{blue}{\textbf{\textit{"For each defined temperature (T), we calculate the average temperature anomalies (MH-minus-PI T anomalies) and compare the results with temperature reconstructions based on ice core $\delta^{18}$O$_p$  values and observed modern spatial $\delta^{18}$O$_p$ -temperature slopes. As shown in Fig. 9a, there is a lack of consensus between the simulated MAT anomalies and the values obtained from proxy records in Greenland. However, it is evident that the MTWA and T-monsoon exhibit a relatively stronger agreement with the observed data (Fig. 9b,c). Table 3 indicates that the proxy-based mean temperature anomaly over Greenland is 1.01 °C. Notably, the MTWA0.5 and sig. MTWA definitions provide a mean temperature anomaly of 1.29 °C across the Greenland ice core locations, which is relatively closer to the observed data compared to other T definitions. In Antarctica, non-uniform temperature anomalies are observed, with modest changes in some areas and more significant cooling and warming in other regions. However, our model shows a general cooling trend across the entire Antarctic continent for all T definitions (Fig. 9d-f)."}}}



18. L373 Again I wonder how the results would look if one based the maps on significant anomalies instead of a somewhat arbitrary threshold? Did you look at the slope for the NH/SH monsoon season (JJAS/DJFM) instead of annual mean?

\medskip

\textcolor{blue}{\textbf{Regarding the first question, we refer to our answer to the 16th comment above. }}

\medskip

\textcolor{blue}{\textbf{Regarding the second question, yes we calculated the isotope-temperature gradient for the monsoon seasons (JJAS for the Northern Hemisphere, DJFM for the Southern Hemisphere) as shown in Fig.~\ref{temporal}c of this response-letter. There are some unrealistic high values found for the Arctic Ocean and unrealistic negative values obtained for the North America, Bering Sea and mid-latitude Asia. The gradient values for the ice core locations are as follows: Greenland: 0.61 ‰/$^{\circ}C$; Antarctic: 0.28 ‰/$^{\circ}C$; East Antarctica: 0.34 ‰/$^{\circ}C$; West Antarctica: 0.20 ‰/$^{\circ}C$. Thus, the isotope-temperature gradients for the Antarctic areas are too low. Further applying a threshold for the temperature changes (0.5 C) can improve the result (Fig.~\ref{temporal}d), but the isotope-temperature gradients for the Antarctica are still much smaller than observations: Greenland: 0.54 ‰/$^{\circ}C$; Antarctic: 0.32 ‰/$^{\circ}C$; East Antarctica: 0.37 ‰/$^{\circ}C$; West Antarctica: 0.23 ‰/$^{\circ}C$.  If temperature threshold is 1 C, then most areas are excluded from calculation (Fig.~\ref{temporal}e).}}

\medskip
\textcolor{blue}{\textbf{According to the comment, we now added Fig.~\ref{temporal} as Fig. S2 in our supplementary, and we also add the following texts in the revised manuscript (see L405-419 in the difference-tracked version):}}

\textcolor{blue}{\textbf{\textit{"Instead of using an arbitrary  threshold, we propose an alternative method to calculate the temporal gradient of  $\delta^{18}$O$_p$-temperature based on the significance of temperature anomalies. In Fig. S2a, we only consider regions with significant changes in MAT (sig. MAT). Positive gradients are observed in parts of Northern Hemisphere continents; however, data is missing for  Greenland due to the minor change in annual mean temperature between MH and PI, which is statistically insignificant. The simulated temporal isotope-temperature gradient is about 0.6  ‰/$^{\circ}C$ for the Antarctica. Especially, the gradient value of 0.61 ‰/$^{\circ}C$ for West Antarctica is more close to observation compared to other T definitions. For regions with significant changes in MTWA (sig. MTWA), we obtain a gradient pattern similar to that of MTWA0.5 (cf. Fig. S2 and Fig. 8c, and refer to Table 2 for detailed values). Additionally, we calculate the isotope-temperature gradient for the monsoon seasons (sig. T-monsoon), specifically JJAS for the Northern Hemisphere and DJFM for the Southern Hemisphere.  As depicted in Fig. S2c,  unrealistic high values are observed in the Arctic Ocean, while unrealistic negative values are obtained for North America, the Bering Sea, and mid-latitude Asia. The MH-to-PI temporal isotope-temperature gradient for Greenland ice core locations is estimated to be 0.61 ‰/$^{\circ}C$. However, the $\delta^{18}$O$_p$-temperature temporal gradient for Antarctica is largely underestimated (Table 2).Applying a threshold of 0.5 °C for temperature changes (sig. T-monsoon0.5) improves the gradient values for Antarctica, but they still remain considerably smaller than the observations  (Fig. S2d, Table 2).  Further increasing the  temperature threshold to 1 °C (sig. T-monsoon1) excludes most areas from our calculation (Fig. S2e)."}}}

\medskip

\textcolor{blue}{\textbf{Furthermore, we also added the new gradient numbers in Table 2 in the revised paper.}}

\medskip


19. Section 6: In relation to this section the definition of the calendar. Maybe it’s because the calendar effect is neglectable given the variance of the onset of the monsoon, but the calendar effect is a constant bias, so I suggest taking it into account. In any case I think it would be useful to discuss the role of the calendar effect in relation to the onset and duration of the monsoon.

\medskip

\textcolor{blue}{\textbf{It is important to consider the calendar effect if we calculate monthly or seasonal mean values. Given that the analysis on WASM onset is based on only daily data, we can stick to the modern classical calendar. However, if we use the angular calendar, it is important to describe the onset as "day of the year" (e.g., the 174th day of the year) instead of using the specific date (e.g., June 20th), because the lengths of months are different for PI and MH, therefore using the specific date is sometimes misleading. Moreover, if we use  "day of the year" as the unit of WASM onset, it is very simple to convert the onset values from classical calendar to angular calendar, because there is just a constant shift between the two calendars: the first day of the year for the MH is 3 days advanced in the angular calendar (corresponding to 29th December in classical calendar). Therefore, if the onset happens on the 174th day of the year in classical calendar, then it should be on the 177th day of the year in angular calendar.}}

\medskip
\textcolor{blue}{\textbf{According to the comment, we now added in the end of section 6 (see also L541-544 in the difference-tracked version):}}

\textcolor{blue}{\textbf{\textit{"Moreover, as the lengths of months and seasons are different between MH and PI angular calendars \citep{shi2022calendar},  the date of WASM onset in this section is based only on modern classical calendar. We keep in mind that there is a constant shift between the two types of calendars by 1-3 days for PI and MH \citep{shi2022calendar},  however, this effect is neglectable given the large variance of the onset of the WASM."}}}


\medskip

20. L483 From the model setup I understood that the model is run with active vegetation? L141-142: “In our simulations, the dynamic vegetation is interactively calculated via the land surface model JSBACH.” If not, lacking vegetation-albedo feedbacks could explain some of the too weak modelled climate response to MH conditions, if the vegetation is indeed interactive the model might still be underestimating the changes in vegetation.

\medskip

\textcolor{blue}{\textbf{Yes the model includes dynamic vegetation, to make it clearer, we now wrote  (L160 in the difference-tracked version):}}

\medskip

\textcolor{blue}{\textbf{\textit{"In our simulations, the dynamic vegetation is calculated via the land surface model JSBACH."}}}

\medskip

\textcolor{blue}{\textbf{and}}

\medskip

\textcolor{blue}{\textbf{\textit{"The atmospheric component of the model is ECHAM6 \citep{stevens2013atmospheric} which also contains  a land-surface module (JSBACH)  representing multiple plant functional types and two types of bare surface \citep{loveland2000development,raddatz2007will}."}}}
 
\medskip

References

Cuffey, K., R. Alley, P. Grootes, and S. Anandakrishnan (1992), Toward using borehole temperatures to calibrate an isotopic paleothermometer in central Greenland, Global Planet. Change,  98(2-4),  265– 268.

Guan, J.,  Liu, Z.,  Wen, X.,  Brady, E.,  Noone, D.,  Zhu, J., and  Han, J. (2016),  Understanding the temporal slope of the temperature-water isotope relation during the deglaciation using isoCAM3: The slope equation, J. Geophys. Res. Atmos.,  121,  10,342– 10,354, doi:10.1002/2016JD024955.

Kindler, P., Guillevic, M., Baumgartner, M., Schwander, J., Landais, A., and Leuenberger, M.: Temperature reconstruction from 10 to 120 kyr b2k from the NGRIP ice core, Clim. Past, 10, 887–902, https://doi.org/10.5194/cp-10-887-2014, 2014.

Masson-Delmotte, V., et al. (2008), A review of Antarctic surface snow isotopic composition: Observations, atmospheric circulation, and isotopic modeling, J. Clim.,  21(13),  3359– 3387, doi:10.1175/2007JCLI2139.1.

Sime, L., Wolff, E., Oliver, K. et al. Evidence for warmer interglacials in East Antarctic ice cores. Nature462, 342–345 (2009). https://doi.org/10.1038/nature08564

Sjolte, J.,  Hoffmann, G.,  Johnsen, S. J.,  Vinther, B. M.,  Masson-Delmotte, V., and  Sturm, C. (2011), Modeling the water isotopes in Greenland precipitation 1959–2001 with the meso-scale model REMO-iso, J. Geophys. Res.,  116, D18105, doi:10.1029/2010JD015287.

Jesper Sjolte, Georg Hoffmann & Sigfús Jóhann Johnsen (2014) Modelling the response of stable water isotopes in Greenland precipitation to orbital configurations of the previous interglacial, Tellus B: Chemical and Physical Meteorology, 66:1, DOI: 10.3402/tellusb.v66.22872

Vinther, B., Buchardt, S., Clausen, H. et al. Holocene thinning of the Greenland ice sheet. Nature 461, 385–388 (2009). https://doi.org/10.1038/nature08355

Martin Werner, Martin Heimann & Georg Hoffmann (2001) Isotopic composition and origin of polar precipitation in present and glacial climate simulations, Tellus B: Chemical and Physical Meteorology, 53:1,53-71, DOI: 10.3402/tellusb.v53i1.16539



 
 \section{Comments from Reviewer 2}
 
This manuscript presents first results from a stable water isotope enabled Earth system model AWI-ESM-2.1-wiso, provides the information on new implementation of water isotope to the ocean component FSOM. Using the fully coupled AWI-ESM-wiso, two simulations for PI and MH are performed and the model results are evaluated against observed isotope compositions in precipitation and sea water. They conclude that model performs well on capturing the isotopic signals under PI and MH conditions, as well as the changes during the MH compared to PI. The coupled version yields better results than previous version without coupling to the ocean model.
 
The development a fully coupled water isotope enabled ESM will have significant contribution to understand the past climate change and climate variability. The manuscript is well written, documents the comprehensive information about this new development, provides an important reference for future researches based on this model. I have one major comment for authors to consider, may help to improve the manuscript for final publication.

\medskip
 
 \textcolor{blue}{\textbf{We really appreciate your positive comments and suggestions, which have significantly improved our manuscript. We have modified our paper accordingly. Details are as follows: }}

\medskip
 
As an example to show the most distinct changes in WASM during the MH in view of water isotope, the authors propose to use the water isotope to identify the onset of WASM instead of using precipitation, which is regarded as a novel method in this work. The presentation for this part can be improved and make it more convincing. When using the daily data to define the onset of the monsoon, I suppose the authors used 100 years mean daily data for the analysis. Considering high frequency variability in precipitation, which may exceed the climate change between MH and PI, 100-years mean for daily data potentially can remove the high frequency variability and largely smooth the data, consequently yield the no-change in onset date as presented in the results. However, from the frequency showed in Fig12 histogram, one can observe that isotope based method shows about 10-day earlier in monsoon onset in MH than that of PI, which is a reasonable results for the longer NH summer under MH orbital forcing. Using the similar statistic for the monsoon withdraw date, one can estimate the duration of the WASM during MH and PI. I think the focus on the change of duration of monsoon season makes more sense than onset date of monsoon for past monsoon.
 
I suggest authors revise this part by identifying monsoon onset, withdraw and duration for individual model year, and using the statistics (as show in Fig12) for MH and PI comparison. The advantage of using isotope instead of precipitation to define the monsoon season may still hold true, but more convincing.

 \medskip

                \begin{figure}
      \centering\includegraphics[width=0.7\textwidth]{hist_onset.pdf}
        \caption{(a,b) Histogram of WASM onset based on the precipitation approach (solid bars) and isotope approach (stippling bars) for (a) PI and (b) MH. (c,d) Histogram of WASM withdrawal for (c) PI and (d) MH. (e,f) Same as (c,d) but for WASM duration. \label{hist}}
    \end{figure}
    
 \textcolor{blue}{\textbf{Thanks very much for the constructive comment,  the WASM onset was computed based on individual model year, we agree with the reviewer that  the duration of monsoon season shall be explored. For this purpose, we now added analysis on the withdrawl of WASM,  and the monsoon’s duration is simply computed as the time interval between its onset and withdrawal. Moreover, we performed histogram analysis on the WASM onset, withdrawal and duration, which is shown in Fig.~\ref{hist} in this response letter and also in Fig. 12 in the revised manuscript.}}
 
 \textcolor{blue}{\textbf{Here are the related texts added/modified in the manuscript: (please also refer to L498-535 in the difference-tracked version)}}

 \textcolor{blue}{\textbf{\textit{"Analyses of our PI experiment suggest a mean WASM onset date of June 19th  based on the isotope approach with a standard deviation of 14 days,  the MH simulation yields a similar result (June 20th $\pm$ 11 days),  indicating that there is no discernible difference between the MH and PI in terms of the simulated  initiation date  of WASM. The initiation date of WASM, as determined by precipitation alone, is June 21th $\pm$ 15 day for PI, similar to the isotope-based approach. For MH,  the precipitation-based method yields a  WASM onset date of June 23th $\pm$ 12 days.  In addition, our histogram plot (i.e., Fig. 12a,b) shows that the WASM onset defined by both approaches is  within 140-200th day of the year. Our result also reveals that the onset of the WASM occurs most frequently between the 160th and 180th day of the year."}}}

 \textcolor{blue}{\textbf{\textit{"With respect to the withdrawal of the WASM, we exclusively employ the daily precipitation time series.   According to \citet{pausata2016impacts}, we define the WASM withdrawal as the date on which the zonally averaged precipitation at the north latitude of ITCZ falls below 2 mm/day, for at least 20 consecutive days. The monsoon’s duration is simply computed as the time interval between its onset and withdrawal. Based on this approach, the calculated WASM withdrawal is  11th October and September 30th for  PI and MH, with a standard deviation of 13 days and 7 days respectively. Thus the duration of WASM is 114 days for PI and 102 days respectively, indicating a shorter WASM trimester in MH. In PI, the retreat of the summer monsoon over West Africa is most frequently observed within 290-320th day of the year (Fig. 12c,d), while in most model years, the WASM in MH terminates within 280-300th day of the year. The duration of WASM in PI can reach up to 180 days whereas for MH the maximum WASM length is only 150 days (Fig. 12e,f). Furthermore, following  \citet{d2019northern}, we define the WASM domain as the area where the precipitation anomaly between the JJAS and DJFM seasons exceeds 2 mm/day. To examine the retreat of the summer monsoon on a local scale,  we further  calculate the termination date of the WASM for each grid box within the monsoon domain. As depicted in Fig. 13a,b, the termination of the rainy season exhibits a meridional distribution, driven by the shift in maximum insolation. We also observe an earlier withdrawal of the WASM in most parts of the monsoon domain, particularly in the southern region (Fig. 13c). Additionally, the shorter duration of MH WASM can be attributed to a more rapid decrease in insolation at the north location of the ITCZ in MH compared to PI. Our result indicates that the simulated intensification in MH WASM is attributable to  a rise in precipitation rate alone than an extension of the monsoon season. "}}}
 
  \medskip
  
Minor comments:

 \medskip
 
1. Line 80-84, should motivate why studying the monsoon onset is important. Line 495-507 in the discussion part can be moved here.

\medskip

\textcolor{blue}{\textbf{Thanks for the comment, we now modified the text into (see also L93-101 in the difference-tracked version):}}

\textcolor{blue}{\textbf{\textit{"The subtropic rainfall was  isotopically more depleted \citep{herold2009eemian,cauquoin2019water,gierz2017simulating} as a result of enhanced precipitation and stronger advection of moisture from the source ocean surface \citep{herold2009eemian}. A variety of studies have attempted to derive the modern summer monsoon onset based on observations for the region of Asia \citep{nguyen2014climatological,moron2014interannual,joseph2006summer}, Africa  \citep{sultan2003west,fitzpatrick2015west,dunning2016onset}, and North America \citep{bombardi2020detection}. However, because proxy time series cannot resolve daily variations in precipitation or water isotopes,  there is a lack of knowledge on the initiation date of MH summer monsoon due to low temporal resolution in proxy data. Since models can give adequate high-frequency variables to explore synoptic phenomena, we can study the possible characteristics of monsoon onset using the model outputs with high temporal resolution under both PI and MH climatic conditions."}}}

\medskip


2. Line 125, "The model has been widely used with its standard configuration", may cause confusion if it refers to model with or without isotope, better present "The AWI-ESM model has been widely used with its standard configuration".

\medskip

\textcolor{blue}{\textbf{Thanks for the comment, we now added "AWI-ESM" in the sentence (see also L143 in the difference-tracked version).}}

\medskip

3. This is the first document for implementing the isotope in ocean model, is it possible to evaluate if everything is done correctly with ocean-only simulation without coupling? Otherwise how to identify any biases in coupled version is due to the ocean model implementation or due to coupling?

\medskip

\textcolor{blue}{\textbf{Description of atmospheric processes as well as isotope fractionation can surely have an impact on the simulated results, on the other hand,  the ocean dynamics impact the distribution of passive tracers (including water isotopes) which in turn impact the resulting isotopic concentrations in precipitation, as ocean is the main source of moisture. Therefore, any unrealistic treatment in the atmosphere or ocean model can lead to bias in simulated isotopic values. Currently we do not have simulations performed using the standard-alone ocean model. But there are simulations by ECHAM6-wiso. However, it is still a challenge to estimate whether the bias origins from the air or the ocean, as for running the stand-alone atmosphere model we need to provide oceanic boundary conditions like the sea surface temperature, sea ice concentration and the isotope composition of surface sea water, and the simulated results depend significantly on these boundary conditions. For example, using AWIESM-wiso, the simulated LGM has a warming bias over Greenland and this leads to an overestimation of LGM $\delta^{18}$O$_p$ over that region. But, using ECHAM6-wiso forced by reconstructed SST and GLOMAP sea ice, the realistic ocean conditions provide a good constraint on the simulated LGM climate and the isotope values. Thus, it is hard to say how many percent of this improvement is due to the boundary conditions or the removal of ocean module.}}

\medskip


4. A bit confusing for the PI period, Line 130 says 1850 CE, and Line 177 says "and PI period (1950-1990CE)".

\medskip

\textcolor{blue}{\textbf{Sorry for the typo, for speleothem data we need to select a time range to represent PI and to be compared to the simulated PI values, and we haven chosen the time period of 1850-1990CE. In the revised version we corrected this information (see also L196 in the difference-tracked version):}}

\textcolor{blue}{\textbf{\textit{"Following  \citet{comas2019evaluating}, in the present study we select 30 speleothem sites (33 cores) for which isotopic values are available for both the MH (defined as the time interval of 6 ± 0.5 ka) and modern periods (1850–1990 CE)."}}}

\medskip

5. Line 226-227, "For instance, our model underestimates the isotopic composition of the South Atlantic and overestimates the δ18O values in the Southern Ocean", these under/over estimation is due to ocean model or due to implemented isotope? Any comments on these biases would be helpful to understand the origin of the biases.

\medskip

\textcolor{blue}{\textbf{Now we added some discussion on the potential reason for this model bias (see also L247-251 in the difference-tracked version):}}

\textcolor{blue}{\textbf{\textit{"For instance, our model underestimates the  isotopic composition  of the South Atlantic and overestimates the $\delta^{18}$O values in the Southern Ocean. Mismatches between our model simulation and marine $\delta^{18}$O reanalysis data may result from model bias in representing freshwater export or net precipitation. Moreover, a warming bias across the Southern Ocean simulated by AWIESM-wiso  can lead to increased evaporation that favors larger $\delta^{18}$O$_{oce}$ values."}}}

\medskip



6. Colour bar in Fig 1,2,5,6,8 needs to be improved. For example, the red colour in Fig.1 for -6 to -4 looks the same to me.

\medskip

\textcolor{blue}{\textbf{In the revised manuscript, we have changed the color table for the mentioned plots to make the color blocks more distinguishable from each other. Moreover, we also pay attention to that all color tables used in the paper are suitable for people who needs aid in colors.}}

\medskip

7. To be consistent, mark RMSE in Fig6 as that in Fig5.

\medskip

\textcolor{blue}{\textbf{Thanks for the comment, we now added the RMSE across all species in Fig. 6b, and we also display the RMSE value for each individual foraminifera type in the corresponding legend. We refer to   Fig. 6b in our revised manuscript.}}

\medskip

8. Line 374, for gradient, "-4 ‰mm−1d" should be "-4 ‰mm−1/d".

\medskip


\textcolor{blue}{\textbf{Isotope values has the unit of ‰, precipitation has the unit of mm/day, thus the unit for isotope-precipitation gradient is ‰/(mm/day) and is therefore ‰mm−1day. We have correct the value from "-4 ‰mm−1d" to "-4 ‰mm−1day" in order to be consistent with the precipitation unit (mm/day).}}

\medskip


9. Line 387, dexp is never mentioned earlier, for those who are not familiar with water isotopes, the meaning of d-excess and definition of dexp should be introduced.

\medskip


\textcolor{blue}{\textbf{Thanks for the comment, we now modified the text into (see also L446-448 in the difference-tracked version):}}

\textcolor{blue}{\textbf{\textit{"In this section for the first time we examine the onset of mid-Holocene West Africa summer monsoon (WASM)  using the simulated monsoon-related changes of both climatic and isotopic variables (i.e., precipitation,  $\delta^{18}$O$_p$ and deuterium excess) on a daily scale. The deuterium excess (hereafter referred to as $dex$) is a second-order isotopic parameter defined as   $dex =  \delta D-8\delta^{18}$O. Changes in $dex$ are often interpreted as a source region effect, namely related to the humidity and temperature conditions at the evaporation source regions  \citep{merlivat1979global}."}}}

\medskip

10. Fig10, better to indicate "day of a year" for x-axis in the figure, the same for the other similar figs.


\medskip


\textcolor{blue}{\textbf{Thanks for the suggestion, we have added the information for x-axis in the corresponding figures.}}


\medskip

11. Fig10 caption, "The two black lines respond to 4.7°N and 10.3°N." There is one black line and one red line. Here the 10.3°N is marked in the figure, but in the text description often mentioned 12.3°N, typo?

\medskip


\textcolor{blue}{\textbf{Sorry for the typo. The number should be 10.3 instead of 12.3, we have corrected it in the revised version. Thanks.}}


\medskip


12. For Fig S2 to S4, is it proper to use 4.7°N and 10.3°N for the specific year? If you apply EOF analysis for specific year, I think the maxima precipitation will have different latitude location.

 \medskip

\textcolor{blue}{\textbf{It is a good idea to check the location of maxima precipitation in these special cases. Here we performed EOF for these model years, and the result is shown in Fig.~\ref{eof} in this response letter (below). From  Fig.~\ref{eof}e we can see that the maxima precipitation from EOF1 are all located at 10.3 °N. }}




\begin{figure}
\centering
 \noindent\includegraphics[width=0.8\textwidth]{step1_year.pdf}
\caption{(a-d) Spatial pattern of EOF1 precipitation over West Africa for the 4 cases in Fig. S2-S5 in the original manuscript (which corresponds to Fig. S3-S6 in the revised manuscript). (e) EOF1 precipitation averaged over each latitude. \label{eof}}
\end{figure}

\medskip

\bibliographystyle{copernicus}
\bibliography{ref.bib}

\end{document}
